
回忆
$E_{((\mathfrak t_1,\ell_1),w_1,\vec p^{\,1})}$
与
$E_{((\mathfrak t_0,\ell_0),w_0,\vec p^{\,0})}$
中的元素均具有紧支集。
以下简记这两个空间为 $\mathcal O_1,\mathcal O_0$，
并把胶合后所选的障碍空间
$E_{((\mathfrak t',\ell'_T),w'_T,\vec p')}$
记为 $\mathcal O_T$。
利用紧支性，可选择 $\mathcal O_T$ 使
\[
J^S_{1,*}+J^S_{0,*}\colon
\mathcal O_1\oplus\mathcal O_0\longrightarrow\mathcal O_T
\]
为单射。因此将 $\mathcal O_1\oplus\mathcal O_0$
视作 $\mathcal O_T$ 的子空间。
现在考虑方程
\begin{equation}
\bar\partial w\equiv0\pmod{\mathcal O_1\oplus\mathcal O_0}.
\tag{7.1.22}\label{eq:7-1-22}
\end{equation}
这里 $w$ 在 $C^0$ 意义下接近 $w'_T$。

先描述式~\eqref{eq:7-1-22} 的近似解。
为此，引入\emph{小切空间}
\[
\begin{aligned}
\mathcal T_1
&:=T_{((\mathfrak t_1,\ell_1),w_1,\vec p^{\,1})}
 W_{k+1}^{1,p}(M,L),\\
T^{\mathrm{small}}_{((\mathfrak t_1,\ell_1),w_1,\vec p^{\,1})}
 W_{k+1}^{1,p}(M,L)
&:=\bigl\{(\xi_1,V_1,\vec v^{\,1})\in\mathcal T_1
\mid\xi_1=0\bigr\}.
\end{aligned}
\]
同样定义
$T^{\mathrm{small}}_{((\mathfrak t_0,\ell_0),w_0,\vec p^{\,0})}
 W_{m+1}^{1,p}(M,L)$。
在式~\eqref{eq:7-1-19-1}--\eqref{eq:7-1-19-2}
中以这两个小切空间代替全切空间，
得到 $\mathcal X^+_{\mathrm{small}}$
和 $\mathcal X_{\mathrm{small}}$。
这里不用整个 $\mathcal X$，
是因为我们要以不同于 $V_1,v_0^1,\ldots,v_k^1$
等参数的方式，单独处理模参数 $\xi_1,\xi_0$。
\translatornote{原书此处把第二分支 $w_0$ 的空间下标写成 $k+1$，把两个模参数写成 $\xi_1,\xi_2$；但式 (7.1.19.1) 已固定第二分支为 $W_{m+1}^{1,p}$，参数为 $\xi_0$。下文 $I^S$ 的目标下标及输入元组次序也依同一型别统一。见疑误 A-0415。}

对 $S>2$ 定义映射
\[
I^S\colon\mathcal X_{\mathrm{small}}\longrightarrow
T_{((\mathfrak t',\ell'_T),w'_T,\vec p')}
W_{k+m}^{1,p}(M,L).
\]
这里 $\vec p'=(p'_1,\ldots,p'_{k+m})$
由下列点依胶合后的边界次序重新排列而成：
\[
p_1^1,\ldots,p_k^1,
p_0^0,\ldots,\widehat{p_j^0},\ldots,p_m^0.
\]
取
$((0,V_1,\vec v^{\,1}),(0,V_0,\vec v^{\,0}))
\in\mathcal X_{\mathrm{small}}$，并记
$v=v_0^1=v_j^0$，见式~\eqref{eq:7-1-19-2}。
为简写分段式，令
\[
\begin{aligned}
I_{i,T}&:=I_{(\mathfrak t_i,\ell_i),(\mathfrak t',\ell'_T)},
\quad i=0,1,\\
\mathcal D_0&:=I_{0,T}\bigl(X(\mathfrak t_0,\ell_0)
\setminus(( -\infty,5T)\times[0,1])\bigr),\\
\mathcal D_1&:=I_{1,T}\bigl(X(\mathfrak t_1,\ell_1)
\setminus((-5T,\infty)\times[0,1])\bigr),\\
B_S(\tau,t)&:=v+\chi^+_{-S}(\tau)
 \bigl(\operatorname{Pal}(V_0(\tau,t))-v\bigr)
 +\chi^-_{+S}(\tau)
 \bigl(\operatorname{Pal}(V_1(\tau,t))-v\bigr).
\end{aligned}
\]
其中 $\operatorname{Pal}$ 均表示沿相应最短测地线的平行移动。
于是 $I^S$ 由下式给出：
\begin{equation}
\begin{aligned}
&I^S\bigl((0,V_1,\vec v^{\,1}),(0,V_0,\vec v^{\,0})\bigr)(x)\\
&\qquad=
\begin{cases}
V_0(x),&x\in\mathcal D_0,\\
\operatorname{Pal}\bigl(B_S(\tau,t)\bigr),
 &x=(\tau,t)\in[-5T,5T]\times[0,1],\\
V_1(x),&x\in\mathcal D_1.
\end{cases}
\end{aligned}
\tag{7.1.23}\label{eq:7-1-23}
\end{equation}
此处输入依 $\mathcal X_{\mathrm{small}}$
的定义先取分支 $1$，再取分支 $0$；
原书在式~(7.1.23) 把两项写成反序，
而式~(7.1.24) 又恢复了此处所用的次序。

将 $I^T$ 限制到相应核空间，得到
\translatornote{原书将这个核写成 $D_{w'_T}\bar\partial$ 在 $\mathcal X_{\mathrm{small}}$ 上的限制；$\mathcal X_{\mathrm{small}}$ 属断开的两分支，故依式 (7.1.19) 与 (7.1.26) 改为 $D_{(w_1,w_0)}\bar\partial$。见疑误 A-0414。}
\[
I^T\colon
\operatorname{Ker}\bigl(
\pi\circ D_{(w_1,w_0)}\bar\partial
\vert_{\mathcal X_{\mathrm{small}}}\bigr)
\longrightarrow
T^{\mathrm{small}}_{((\mathfrak t',\ell'_T),w'_T,\vec p')}
W_{k+m}^{1,p}(M,L).
\]
取一个到其像的微分同胚
\[
E=(E_1,E_2)\colon
\{(x,v)\in TM\mid |v|<\epsilon\}\longrightarrow M\times M,
\]
满足
\[
E_1(x,v)=x,\qquad
\left.\frac{dE_2(x,tv)}{dt}\right|_{t=0}=v,
\qquad
E(x,v)\in L\times L\quad(x\in L,\ v\in T_xL).
\]
例如，可取使 $L$ 成为全测地子流形的度量 $g$，
由其列维--奇维塔联络的指数映射来给出 $E$。
用 $E$ 定义
\[
\operatorname{Exp}\colon
T^{\mathrm{small}}_{((\mathfrak t',\ell'_T),w'_T,\vec p')}
W_{k+m}^{1,p}(M,L)
\longrightarrow W_{k+m}^{1,p}(M,L),
\]
其中
\[
\operatorname{Exp}(0,V,\vec v)(x)
=E_2\bigl(w'_T(x),V(x)\bigr),
\qquad x\in X(\mathfrak t',\ell'_T).
\]
更准确地说，它只在 $(V,\vec v)$ 的零邻域内定义。
对各
$(\mathfrak t'_0,\ell'_0)$、
$(\mathfrak t'_1,\ell'_1)$、$T$
及
$((0,V_1,\vec v^{\,1}),(0,V_0,\vec v^{\,0}))
\in\mathcal X_{\mathrm{small}}$，
我们取
\begin{equation}
\bigl((\mathfrak t',\ell'_T),
\operatorname{Exp}(I^T((0,V_1,\vec v^{\,1}),
(0,V_0,\vec v^{\,0})))\bigr)
\tag{7.1.24}\label{eq:7-1-24}
\end{equation}
作为近似解；其中
$(\mathfrak t',\ell'_T)
=\operatorname{Comp}((\mathfrak t'_1,\ell'_1),
(\mathfrak t'_0,\ell'_0);T)$。

下一步是扰动式~\eqref{eq:7-1-24}，
使之成为式~\eqref{eq:7-1-22} 的真解。
如前所述，为此采用稍非标准的加权索伯列夫范数。
所用权函数是
\[
\alpha\colon X(\mathfrak t',\ell'_T)\longrightarrow[1,\infty).
\]
回忆 $|\cdot|'\colon\mathbb R\to[0,\infty)$
为光滑函数，在 $|x|\geq1$ 时满足 $|x|'=|x|$，
在 $|x|<1/2$ 时满足 $|x|'=0$。
每条外边 $e\in C^1_{\mathrm{ext}}(\mathfrak t')$
对应 $L_e\cong(-\infty,0]\times[0,1]$，
在其上令
\[
\alpha(\tau,t)=\exp(\delta|\tau|').
\]
对于新边 $e(\mathrm{new})$，
将 $L_{e(\mathrm{new})}$ 识别为
$[-5T,5T]\times[0,1]$，并令
\[
\alpha(\tau,t)=
\begin{cases}
1,&|\tau|>5T-1,\\
\exp\bigl(\delta(5T-|\tau|)\bigr),
 &5T-2>|\tau|>1.
\end{cases}
\]
在中间的过渡区域，存在常数 $C$ 使
\translatornote{原书在 $|\tau|<1$ 给出绝对差界 $|\alpha-e^{5T\delta}|<C$；若 $C$ 应与 $T$ 无关，这与相邻区间的指数权重在 $|\tau|\downarrow1$ 时不能连续衔接。原文未说明该处常数是否依赖 $T$，故保留其陈述并标明疑点。见疑误 A-0419。}
\[
1\leq\alpha(\tau,t)<C
\quad(5T-1>|\tau|>5T-2),
\qquad
|\alpha(\tau,t)-e^{5T\delta}|<C
\quad(|\tau|<1).
\]
其他 $L_e$ 上令 $\alpha=1$，见图~\ref{fig:7-1-5}。

\par\medskip\noindent\begin{minipage}{\textwidth}
\centering
\setlength{\unitlength}{0.68mm}
\begin{picture}(145,72)
\linethickness{0.45pt}
\put(8,12){\vector(1,0){130}}
\put(73,3){\vector(0,1){62}}
\put(12,20){\line(1,0){15}}
\qbezier(27,20)(45,20)(59,37)
\qbezier(59,37)(67,47)(73,47)
\qbezier(73,47)(79,47)(87,37)
\qbezier(87,37)(101,20)(119,20)
\put(119,20){\line(1,0){15}}
\put(64,12){\line(0,1){20}}
\put(82,12){\line(0,1){20}}
\qbezier[12](52,63)(61,39)(73,37)
\qbezier[12](73,37)(85,39)(94,63)
\put(4,60){$\scriptstyle e^{\delta(5T-\tau)}$}
\put(104,60){$\scriptstyle e^{\delta(5T+\tau)}$}
\put(13,4){$\scriptstyle -5T+1$}
\put(59,4){$\scriptstyle -1$}
\put(74,4){$\scriptstyle 0$}
\put(84,4){$\scriptstyle 1$}
\put(110,4){$\scriptstyle 5T-1$}
\put(76,20){$\scriptstyle 1$}
\put(138,10){$\scriptstyle\tau$}
\end{picture}
\par\smallskip
{\renewcommand{\thefigure}{7.1.5}
\refstepcounter{figure}
{\normalfont\bfseries 图~\thefigure}\label{fig:7-1-5}}
\quad$\alpha(\tau,t)$
\end{minipage}\par\medskip
下面用 $\alpha$ 定义所需的索伯列夫范数。
设
$w:(X(\mathfrak t',\ell'_T),
\partial X(\mathfrak t',\ell'_T))\to(M,L)$
在 $C^0$ 意义下接近 $w'_T$。
在向量空间
$\mathcal E^{0,p}_{((\mathfrak t',\ell'_T),w,\vec p')}$
上改赋另一范数：对其中的 $V$，令
\begin{equation}
\|V\|_{0,p,\alpha}^p
=\int_{X(\mathfrak t',\ell'_T)}\alpha|V|^p.
\tag{7.1.25}\label{eq:7-1-25}
\end{equation}
在新边 $L_{e(\mathrm{new})}$ 以外，
该权函数同定义 $\mathcal E$ 时所用的权函数一致，
见式~\eqref{eq:7-1-7-3}。
因此式~\eqref{eq:7-1-25} 的右端有限，
对固定的 $T$，$\|\cdot\|_{0,p,\alpha}$
与原范数 $\|\cdot\|_{0,p}$ 等价；
不过两范数的比较比率随 $T\to\infty$ 而无界。

还要改变胶合后小切空间上的范数。
设 $(0,V,\vec v)=(V,\vec v)$
属于下列小切空间：
\[
T^{\mathrm{small}}_{((\mathfrak t',\ell'_T),w,\vec p')}
W_{k+m}^{1,p}(M,L).
\]
当 $p$ 充分大时 $V$ 连续，
因而 $V(0,1/2)$ 有意义；这里
$(0,1/2)\in L_{e(\mathrm{new})}
\cong[-5T,5T]\times[0,1]$。
为使长式清楚，记
\[
V_i^\circ:=V-\operatorname{Pal}(v_i),
\qquad
V_{\mathrm{new}}^\circ
:=V-\operatorname{Pal}(V(0,1/2)).
\]
定义
\[
\begin{aligned}
\|(V,\vec v)\|_{1,p,\alpha}^p
={}&\sum_{\substack{e\in C^1_{\mathrm{int}}(\mathfrak t')\\
                       e\ne e(\mathrm{new})}}
\int_{L_e}(|\nabla V|+|V|)^p\\
&+\sum_{e_i\in C^1_{\mathrm{ext}}(\mathfrak t')}
\int_{L_{e_i}}\alpha
\bigl(|\nabla V_i^\circ|^p+|V_i^\circ|^p\bigr)\\
&+\int_{L_{e(\mathrm{new})}}\alpha
\bigl(|\nabla V_{\mathrm{new}}^\circ|^p
      +|V_{\mathrm{new}}^\circ|^p\bigr)\\
&+\sum_i|v_i|^p+|V(0,1/2)|^p.
\end{aligned}
\]
这里 $\operatorname{Pal}$ 是沿最短测地线的适当平行移动。
\translatornote{原书两条带权积分中的 $p$ 次方括号错置，若按印刷读法会把梯度项变成 $p^2$ 次；译文依同节式 (7.1.6.5) 的 $W^{1,p}$ 范数形式校正，末项还删除一个游离右括号。见疑误 A-0418。}

与引理~\ref{lem:7-1-20} 的证明一样，
可假设 $\pi\circ D_{(w_1,w_0)}\bar\partial$
在 $\mathcal X_{\mathrm{small}}$ 上的限制
\begin{equation}
\pi\circ D_{(w_1,w_0)}\bar\partial\colon
\mathcal X_{\mathrm{small}}
\longrightarrow
\frac{\mathcal E^{0,p}_{((\mathfrak t_1,\ell_1),w_1,\vec p^{\,1})}}
{\mathcal O_1}
\oplus
\frac{\mathcal E^{0,p}_{((\mathfrak t_0,\ell_0),w_0,\vec p^{\,0})}}
{\mathcal O_0}
\tag{7.1.26}\label{eq:7-1-26}
\end{equation}
为满射。

\begin{proposition}{7.1.27}\label{prop:7-1-27}
假设式~\eqref{eq:7-1-26} 为满射。
则存在与 $T$ 无关的常数 $C>0$，具有如下性质：
对每个
$(\mathfrak t',\ell'_T)
=\operatorname{Comp}((\mathfrak t'_1,\ell'_1),
(\mathfrak t'_0,\ell'_0);T)$
和每个充分接近 $w'_T$ 的 $w$，
存在线性映射
\[
\begin{aligned}
Q_{((\mathfrak t',\ell'_T),w)}\colon
&\frac{\mathcal E^{0,p}_{((\mathfrak t',\ell'_T),w,\vec p')}}
{\mathcal O_1\oplus\mathcal O_0}\\
&\longrightarrow
T^{\mathrm{small}}_{((\mathfrak t',\ell'_T),w,\vec p')}
W_{k+m}^{1,p}(M,L),
\end{aligned}
\]
满足
\begin{align}
(\pi\circ D_w\bar\partial)
\circ Q_{((\mathfrak t',\ell'_T),w)}
&=\operatorname{id};
\tag{7.1.28.1}\label{eq:7-1-28-1}\\
\bigl\|Q_{((\mathfrak t',\ell'_T),w)}(V)\bigr\|_{1,p,\alpha}
&<C\|V\|_{0,p,\alpha}.
\tag{7.1.28.2}\label{eq:7-1-28-2}
\end{align}
\end{proposition}

要强调的是，若以不带权的
$\|\cdot\|_{1,p},\|\cdot\|_{0,p}$
代替加权范数
$\|\cdot\|_{1,p,\alpha},\|\cdot\|_{0,p,\alpha}$，
就不能使常数 $C$ 与 $T$ 无关。
在证明命题~\ref{prop:7-1-27} 以前，
先给出下面的引理；其证明留待本小节末尾。
\translatornote{原书引理右端写 $\|V_2\|_{1,p}$，但输入只有 $V_1,V_0$ 两个分支；译文依输入改为 $\|V_0\|_{1,p}$。见疑误 A-0416。}

\begin{lemma}{7.1.29}\label{lem:7-1-29}
存在与 $T$ 无关的常数 $c>0$、$C>0$，使得
\[
\begin{aligned}
&\bigl\|\bar\partial\operatorname{Exp}
\bigl(I^T((0,V_1,\vec v^{\,1}),
(0,V_0,\vec v^{\,0}))\bigr)\bigr\|_{0,p,\alpha}\\
&\qquad<C\bigl(\|V_1\|_{1,p}
 +\|V_0\|_{1,p}+1\bigr)e^{-cT}.
\end{aligned}
\]
\end{lemma}

一旦得到命题~\ref{prop:7-1-27}
和引理~\ref{lem:7-1-29}，
便可用标准方法构造式~\eqref{eq:7-1-22} 的一族解，
其参数空间是
$\pi\circ D_{(w_1,w_0)}\bar\partial$ 的核
与 $(\mathfrak t,\ell)$ 在
$\overline{\operatorname{Gr}}_{k+m}$
中的一个邻域的直积。
因此，要为无穷远处的点（即低维层中的点）
构造仓西邻域，
只剩证明命题~\ref{prop:7-1-27}
和引理~\ref{lem:7-1-29}。
